\documentclass[11pt]{article}
\usepackage{../homework}

\profname{Dr. Stanislaw Szarek}
\classcode{MATH 423 Measure Theory}
\assignnum{Assignment 09}
\duedate{02 November 2026}

\begin{document}
All problems from Folland, \textit{Real Analysis} (2e).

\begin{problem}{2.3.27}
  Let \( f_n(x) = ae^{-nax} - be^{-nbx} \) where \( 0 < a < b \).
  \begin{enumerate}[label=(\alph*)]
    \item \( \sum\limits_1^{\infty} \int\limits_0^{\infty} |f_n(x)|\dd{x} = \infty \).
          \begin{proof}
            Since \( |f_n| \) is non-negative and continuous, it is measurable and in \( L^+ \), so that we may interchange the sum and the integral.
            Then,
            \begin{equation}
            \label{eq:2}
            \int\limits_0^{\infty} \sum\limits_{1}^{\infty} | ae^{-nax} - be^{-nbx} | \dd{x},
          \end{equation}
          \end{proof}
    \item \( \sum\limits_1^{\infty}\int\limits_0^{\infty} f_n(x)\dd{x}  = 0 \).
    \item \( \sum\limits_1^{\infty} f_n \in L^1([0,\infty], m) \), and \( \int\limits_0^{\infty} \sum\limits_1^{\infty} f_n(x) \dd{x} = \log(b/a) \).
  \end{enumerate}
\end{problem}

\begin{problem}{2.3.31(a)}
  Derive the formula by expanding part of the integrand into an infinite series and justifying th term-byterm integration. Exercise 29 may be useful.
  \begin{enumerate}[label=(\alph*)]
          \item For \( a>0 \), \( \int\limits_{-\infty}^{\infty} e^{-x^2}\cos ax\dd{x} = \sqrt{\pi} e^{-a^2/4} \).
  \end{enumerate}
\end{problem}

\begin{problem}{2.4.39}
  If \( f_n \rightarrow f \) almost uniformly, then \( f_n \rightarrow f \) a.e. and in measure.
\end{problem}
\begin{proof}
  By definition, \( f_n \rightarrow f \) almost uniformly means that for every \(\epsilon > 0\) there is \( E \subset X \) with \(\mu(E) < \epsilon\) and \( f_n \rightarrow f \) uniformly on \( X \setminus E \).
  In particular, for all \( n \in \mathbb{N} \) there is \( E_n \) with \(\mu(E_n) < 1/n\), whence we define
  \begin{equation}
  \label{eq:1}
  E = \bigcap_n E_n.
\end{equation}
By continuity from above (\( \mu(E_1) < 1 \)),
\begin{equation}
\label{eq:3}
\mu(E) = \lim \mu(E_n) = 0.
\end{equation}
\end{proof}

\begin{problem}{2.4.42}
  Let \(\mu\) be the counting measure on \( \mathbb{N} \).
  Then \( f_n \rightarrow f \) in measure iff \( f_n \rightarrow f \) uniformly.
\end{problem}

\end{document}
